\section{Discussion}\label{sec:discussion}

\paragraph{Four lessons, as budget statements.}
\begin{enumerate}[label=(\alph*),leftmargin=*]
  \item With a perfect model, the deficit of classical DA is restriction and
    approximation. It is honest --- the scheme stays calibrated for what it uses
    --- and it is reducible by smoothing, members and iterations.
  \item Under model error, misspecification dominates. It shows in the
    large-$\tau$ score before it shows in the RMSE, and a scalar inflation hides
    it in aggregate without removing it.
  \item A better prior is not a better estimate. An approximate learned prior
    can beat the true equations, because the ordering is set by the
    approximation budget of the scheme that uses them.
  \item Model-based DA can misreport the information value of its conditioning
    variables: with a wrong model, additional observations stop helping, and the
    reported uncertainty does not say so.
\end{enumerate}

\paragraph{What classical DA still gets right.}
It wins in the sparse regime (at most about ten observation times per window).
It transfers across observing systems, because it keeps $\Hop$ and $\mathbf{R}$
explicit. It is calibrated when its assumptions hold, and it needs no training
set. The learned schemes we compare were trained on 1000 windows of the true
system; that data is a resource classical DA does not use, and
Sect.~\ref{sec:q2-oracle} is designed to price it.

\paragraph{Back to the puzzle.}
\todo{What the budgets suggest for reanalysis vs observation-only products:
misspecification under model error, visible as a floor in the marginal value of
observations; and for the skill of learned priors: approximation and
restriction terms of the model-based scheme at a perfect model. Stated as
hypotheses for real systems.}

\paragraph{Limitations.}
\begin{itemize}[leftmargin=*,nosep]
  \item $\mathrm{mmse}_p(\tau)$ is unknown, so the irreducible term is only
    upper-bounded and every share is a share of an excess.
  \item The DA ensembles are scored in Gaussian form; the kernel density is a
    check.
  \item The cross term between misspecification and approximation is bounded by
    contrasts, not identified.
  \item One idealised system: transfer to a quasi-geostrophic system is
    underway.
  \item The learned \So{} rows are oracle bounds until the biased-model arms are
    in.
\end{itemize}

\paragraph{Positioning.}
\todo{Model-error DA: weak-constraint 4D-Var~\cite{tremolet2006weak}, inflation and
model-error estimation~\cite{carrassi2018da,farchi2021model}, unstable-subspace
DA~\cite{trevisan2004unstable,bocquet2017unstable}. Learned and hybrid DA:
\cite{brajard2020combining,bocquet2020bayesian,fablet4dvarnet}; score-based and
flow-based DA~\cite{rozet2023sda}. Proper scores and their decompositions;
mismatched estimation. Verify every entry of \texttt{refs.bib} (all currently
unverified).} The gap we address: no controlled study measures \emph{where} the
error of a DA scheme comes from --- restriction, misspecification or
approximation --- for the mean and for the reported uncertainty, at matched
information across classical and learned schemes.

\paragraph{Outlook.}
The operator~\eqref{eq:psi} suggests a more general programme: writing every DA
algorithm as a flow, and composing the mean of one with the anomaly of another,
as the hybrid does. We leave it to future work, together with self-supervised
adaptation of learned schemes to the observed system and the transfer to a
quasi-geostrophic system.

\section{Conclusion}\label{sec:conclusion}
A restriction costs variance honestly and preserves calibration; a
misspecification costs bias silently and destroys it. On a two-scale Lorenz-96
system, with five classical and five learned schemes under matched information:
\begin{itemize}[leftmargin=*,nosep]
  \item \textbf{Perfect model.} The deficit of classical DA is mostly honest:
    smoothing removes 16--19\,\% of the RMSE and 100 members remove 38\,\% of the
    smoother's MSE.
  \item \textbf{Model error.} Misspecification dominates and shows in the
    large-$\tau$ score (smoother calibration loss 34 $\to$ 63\,\%). The marginal
    value of observations of strong-constraint 4D-Var collapses sevenfold, and a
    weak constraint recovers a quarter of its error.
  \item \textbf{Form of the prior.} The learned prior beats the 30-member smoother
    and not the 100-member one, so the approximation term, not the fidelity of
    the prior, decides.
\end{itemize}
The flow-matching score makes the distinction measurable, for the mean and for
the uncertainty alike.
